\documentclass[11pt]{article}
\usepackage{amsmath,amssymb,amsthm}
\usepackage[margin=1.2in]{geometry}
\newtheorem{theorem}{Theorem}
\newtheorem{lemma}[theorem]{Lemma}
\newtheorem{proposition}[theorem]{Proposition}
\newtheorem{remark}[theorem]{Remark}
\DeclareMathOperator{\ord}{ord}
\newcommand{\fl}[1]{\left\lfloor #1\right\rfloor}
\newcommand{\fpart}[1]{\left\{#1\right\}}

\title{The sum of reciprocals of prime divisors of $\binom{2n}{n}$:\\
the asymptotic for every $n$}
\author{Samuel Lavery}
\date{Draft --- \today}

\begin{document}
\maketitle

\begin{abstract}
Erd\H{o}s, Graham, Ruzsa and Straus (1975) conjectured that
$\sum_{p\mid\binom{2n}{n}}1/p=(1+o(1))\log\log n$.  We prove the conjectured
asymptotic for \emph{every} sufficiently large $n$, executing the route
sketched by Tao (comment on the Erd\H{o}s problem registry, no.~377,
29~August 2025): a band-by-band analysis of the primes
$p\in(n^{1/(k+1)},n^{1/k}]$, in which non-divisibility forces the fractional
parts $\fpart{n/p^j}$ into $[0,\tfrac12)$ at every level $j$, combined with
equidistribution of these fractional parts over primes in each band.  The
combinatorial core of the argument (the Kummer carry criterion, the band
decomposition, and an exact carry-mass conservation law) is machine-verified
in Lean~4.
\end{abstract}

\section{Statement and architecture}

\begin{theorem}[Main]\label{thm:main}
As $n\to\infty$,
\[
\sum_{\substack{p\le n\\ p\mid\binom{2n}{n}}}\frac1p
=(1-o(1))\log\log n .
\]
\end{theorem}

Since $\sum_{p\le n}1/p=\log\log n+M+o(1)$, Theorem~\ref{thm:main} is
equivalent to the statement that the \emph{non-divisor mass}
$E(n):=\sum_{p\le n,\ p\nmid\binom{2n}{n}}1/p$ satisfies $E(n)=o(\log\log n)$
for every large $n$.  (Whether $E(n)=O(1)$ is the full registry problem; it
is not addressed here.)

The proof has four components.

\begin{itemize}
\item[(K)] \textbf{Carry criterion} (Kummer; machine-verified).  For prime
$p$, $p\nmid\binom{2n}{n}$ iff $\fpart{n/p^j}<\tfrac12$ for every $j\ge1$
with $p^j\le 2n$.  [Lean: \texttt{prime\_not\_dvd\_centralBinom\_iff\_carryFree},
file \texttt{Erdos377CarryRails.lean}, standard axiom footprint.]
\item[(B)] \textbf{Band decomposition} (machine-verified).  Splitting primes
by $k=\fl{\log n/\log p}$, the band $(n^{1/(k+1)},n^{1/k}]$ carries Mertens
mass $\log(1+1/k)+O(\cdot)\asymp1/k$, and within one band the level-one
carry condition is a union of explicit ``hyperbola'' intervals.
[Lean: \texttt{levelOneCarryPrimes\_filter\_eq\_sdiff},
\texttt{levelOneCarryMass\_eq\_band\_sum}.]
\item[(E)] \textbf{Band equidistribution} (the analytic core; \S\ref{sec:equi}).
For $k\le(\log n)^{1-\varepsilon}$ and $J=J(k)\to\infty$ slowly, the vector
$\bigl(\fpart{n/p^{k}},\dots,\fpart{n/p^{k-J+1}}\bigr)$ equidistributes as
$p$ ranges over the band, uniformly in $n$; hence the proportion of the
band's mass that is carry-free at the top $J$ levels is
$2^{-J}(1+o(1))$.
\item[(A)] \textbf{Assembly} (\S\ref{sec:assembly}).  Summing
$2^{-J(k)}/k$ over $k\le(\log n)^{1-\varepsilon}$ and treating the range
$k>(\log n)^{1-\varepsilon}$ trivially yields $E(n)=o(\log\log n)$.
\end{itemize}

\section{The exact density law}\label{sec:density}

Within a band, the top-level condition is not uniformly distributed: the
quotient $t=n/p^k$ is small near the band's top edge, and the correct
density is the hyperbola-weighted one.  This is an exact statement, not a
model.

\begin{lemma}[Band measure]\label{lem:measure}
In the coordinate $t=n/p^k$, the Mertens measure $\sum 1/p$ of the band is
log-uniform: for $1\le a<b\le n^{1/(k+1)}$,
\[
\sum_{\substack{p\in(n^{1/(k+1)},\,n^{1/k}]\\ a<n/p^k\le b}}\frac1p
=\frac{\log(b/a)}{k\,\log n}\Bigl(1+O\bigl(E_k(n)\bigr)\Bigr),
\]
and consequently the carry strips $\fpart{t}\in[\tfrac12,1)$ carry the
fraction
$\sigma^{\mathrm c}_k=\bigl(\sum_{m\ge1}\log\tfrac{2m+2}{2m+1}\bigr)/
\log n^{1/(k+1)}$-normalised portion of each dyadic block, tending to
$\tfrac12$ from below at large quotients and biased below $\tfrac12$ at
small quotients.  \emph{[Exact constants to be finalised; the finite band
form is machine-verified as
\texttt{levelOneCarryMass\_eq\_band\_sum} and the two-level form as
\texttt{levelTwoJointCarry\_fiber\_logMass}.]}
\end{lemma}

\begin{remark}[Numerical verification]
For $42$ values of $n\in[10^{11},10^{12}]$ the top-level counts match the
hyperbola-weighted prediction with $z$-scores $-1.04,-1.02,-2.53$ at
$k=2,3,4$ (aggregated); each level below the top halves the count to
within $10^{-3}$ relative.  There is no anomalous ``clock error'': the
only deviations from the naive uniform model are the deterministic
hyperbola weights of this lemma.
\end{remark}

\section{The equidistribution input}\label{sec:equi}

Fix a band index $k$ and write $P=n^{1/k}$.  The required statement is:

\begin{proposition}[Band equidistribution]\label{prop:equi}
Let $\varepsilon>0$, $k\le(\log n)^{1-\varepsilon}$, and $J\le c\log k$.
For any subintervals $I_1,\dots,I_J\subseteq[0,1)$,
\[
\#\bigl\{p\in(n^{1/(k+1)},n^{1/k}]:\ \fpart{n/p^{k-i+1}}\in I_i\ (1\le i\le J)\bigr\}
=\Bigl(\prod_i|I_i|+o(1)\Bigr)\,
\pi\bigl(n^{1/k}\bigr)-\pi\bigl(n^{1/(k+1)}\bigr)\bigr) ,
\]
uniformly in $n$.
\end{proposition}

\noindent\textbf{Reduction to exponential sums.}  By the Erd\H{o}s--Tur\'an--
Koksma inequality in $J$ dimensions, Proposition~\ref{prop:equi} follows from
\begin{equation}\label{eq:weyl}
S(\mathbf h)\;=\;\sum_{p\sim P}
e\!\Bigl(\sum_{i\le J} h_i\,\frac{n}{p^{\,k-i+1}}\Bigr)\;\ll\;
\frac{\pi(P)}{(\log P)^{A}}
\end{equation}
for all $0<\|\mathbf h\|_\infty\le H$ with $H$ a small power of $\log$.

\noindent\textbf{Status of \eqref{eq:weyl}: the active core.}  After
partition into dyadic blocks $p\sim Q$ with
$n^{1/(k+1)}<Q\le n^{1/k}$, the phase has amplitude
$F:=h\,n/Q^{k-i+1}$, and one applies Vaughan's identity, reducing
\eqref{eq:weyl} to two standard families:
\begin{itemize}
\item[(I)] \emph{Type I:}\quad
$\displaystyle\sum_{m\sim M}\Bigl|\sum_{l\sim L}
e\bigl(h\,n/(ml)^{j}\bigr)\Bigr|$,\qquad $ML\asymp Q$, $M\le Q^{1/3}$.
\item[(II)] \emph{Type II (bilinear):}\quad
$\displaystyle\sum_{m\sim M}\sum_{l\sim L}a_m b_l\,
e\bigl(h\,n/(ml)^{j}\bigr)$,\qquad $ML\asymp Q$,
$Q^{1/3}\le M\le Q^{1/2}$, $\|a\|_\infty,\|b\|_\infty\le1$.
\end{itemize}
\textbf{This paper proves both families from first principles; no
analytic result is cited.}  The demand is a saving of a fixed power of
$\log Q$ only (the assembly needs $2^{-J}$ with $J\asymp\log k$), far
below the power-of-$Q$ savings pursued in the exponent-pair literature,
and at this strength each tool reduces to its elementary core:
\begin{itemize}
\item Vaughan's identity: an elementary divisor-sum rearrangement
(\S\ref{sec:vaughan}, self-contained, one page);
\item the Type I bound: Weyl differencing plus the second-derivative
test for the monomial phase $x\mapsto h\,n\,x^{-j}$
(\S\ref{sec:typeI}, self-contained; in carrier language, a phase-cell
cancellation statement);
\item the Type II bound: Cauchy--Schwarz plus an elementary spacing
count for the values $h\,n/(ml)^j$ (\S\ref{sec:typeII},
self-contained).
\end{itemize}
The genuine work, and the reason for Tao's ``lengthy but routine,'' is
uniformity: the amplitude $F=h\,n/Q^{j}$ varies from $\asymp Q$ (top
level, top edge) to large powers of $Q$ (deep quotients), and every
$(F,Q)$ regime occurring for $k\le(\log n)^{1-\varepsilon}$ must be
covered by one of the two self-contained lemmas.  The complete inputs of
this paper are therefore: Kummer's carry criterion and the elementary
central-binomial bounds (both machine-verified), the Chebyshev--Mertens
bounds of \S\ref{sec:density} (machine-verified in finite form), and the
elementary analysis of
\S\S\ref{sec:vaughan}--\ref{sec:typeII}.  Citations are provenance only.

\begin{remark}[No shortcut through PNT alone]\label{rem:noshortcut}
One might hope to avoid \eqref{eq:weyl} for small $k$ by applying the
prime number theorem to each hyperbola component separately.  This fails
for a structural reason worth recording: by Lemma~\ref{lem:measure} the
band's Mertens mass is log-uniform across quotient scales up to
$n^{1/(k+1)}$, while PNT-with-classical-error resolves components only up
to quotient $\exp(c\sqrt{\log P})$ --- a vanishing fraction of the mass in
every band.  The bilinear estimates are therefore needed for \emph{all}
bands, not only deep ones; what varies with $k$ is only the $(F,Q)$
regime.
\end{remark}

\begin{remark}
The band condition at level $k-i+1$ for $i\le J$ involves only the top $J$
levels; the choice $J(k)=\fl{c\log k}$ makes $\sum_k 2^{-J(k)}/k<
\sum_k k^{-1-c\log 2}<\infty$ --- the assembly only needs $J\to\infty$
at logarithmic speed, which is why the equidistribution demand
\eqref{eq:weyl} is mild.
\end{remark}

\section{Elementary oscillatory-sum toolkit}\label{sec:typeI}

All lemmas in this section are classical in substance; they are proved in
full so that the paper is self-contained.  Throughout, $e(x)=e^{2\pi ix}$
and $\|x\|$ is the distance from $x$ to the nearest integer.

\begin{lemma}[Kusmin--Landau]\label{lem:KL}
Let $f\in C^1[a,b]$ with $f'$ monotone and $\|f'(x)\|\ge\delta>0$ on
$[a,b]$.  Then $\bigl|\sum_{a\le l\le b}e(f(l))\bigr|\le 2/(\pi\delta)+1$.
\end{lemma}

\begin{proof}
Write $\theta_l=f(l+1)-f(l)=f'(\xi_l)$ with $\xi_l\in(l,l+1)$; by
monotonicity the $\theta_l$ are monotone and $\|\theta_l\|\ge\delta$.
With $g_l=\bigl(1-e(\theta_l)\bigr)^{-1}$ we have
$e(f(l))=\bigl(e(f(l))-e(f(l+1))\bigr)g_l$, and Abel summation gives
\[
\Bigl|\sum_l e(f(l))\Bigr|
\le 2\max_l|g_l|+\sum_l|g_{l+1}-g_l| .
\]
Since $|1-e(\theta)|=2|\sin\pi\theta|\ge4\|\theta\|$, each
$|g_l|\le(4\delta)^{-1}$.  As $\theta\mapsto(1-e(\theta))^{-1}$ has
bounded variation $O(\delta^{-1})$ on each monotone stretch with
$\|\theta\|\ge\delta$ (the image path is monotone along a circle arc
staying at distance $\ge4\delta$ from $0$), the variation sum is
$O(\delta^{-1})$.  Tracking constants gives the stated bound.
\end{proof}

\begin{lemma}[Second-derivative test]\label{lem:vdc2}
Let $f\in C^2[a,b]$ with $\lambda\le f''\le c\lambda$ on $[a,b]$,
$\lambda>0$.  Then
$\bigl|\sum_{a\le l\le b}e(f(l))\bigr|
\ll_c (b-a)\lambda^{1/2}+\lambda^{-1/2}$.
\end{lemma}

\begin{proof}
$f'$ is increasing with total increase $\le c\lambda(b-a)$.  Fix
$\delta\in(0,\tfrac14)$.  The set where $\|f'\|<\delta$ is a union of at
most $c\lambda(b-a)+2$ intervals, each of length $\le2\delta/\lambda$
(since $f'$ moves at speed $\ge\lambda$); trivially estimate these ranges
by their lengths.  On the complementary intervals apply
Lemma~\ref{lem:KL}.  This gives
$|S|\ll(c\lambda(b-a)+1)\bigl(\delta/\lambda+\delta^{-1}\bigr)$;
choosing $\delta=\lambda^{1/2}$ yields the claim.
\end{proof}

\begin{lemma}[Weyl differencing]\label{lem:weyl}
For any $1\le H\le N$ and $f:[1,N]\to\mathbb R$,
\[
\Bigl|\sum_{l\le N}e(f(l))\Bigr|^2
\le\frac{2N}{H}\sum_{0\le h<H}
\Bigl|\sum_{l\le N-h}e\bigl(f(l+h)-f(l)\bigr)\Bigr| .
\]
\end{lemma}

\begin{proof}
Average the sum over $H$ shifts, apply Cauchy--Schwarz, and expand the
square; the diagonal terms give the $h$-sum.  Standard bookkeeping.
\end{proof}

\begin{lemma}[$r$-th derivative test]\label{lem:vdcr}
Let $r\ge2$ and $f\in C^r[a,b]$ with
$\lambda\le f^{(r)}\le c\lambda$ on $[a,b]$, $N=b-a\ge1$.  Then
\[
\Bigl|\sum_{a\le l\le b}e(f(l))\Bigr|
\ll_{c,r} N\bigl(\lambda N^{r-2}\bigr)^{1/(2^{r}-2)}
+N^{1-2^{2-r}}\lambda^{-1/(2^{r}-2)} .
\]
\end{lemma}

\begin{proof}
Induction on $r$.  The case $r=2$ is Lemma~\ref{lem:vdc2}.  For $r>2$,
apply Lemma~\ref{lem:weyl}; the differenced phase
$f(l+h)-f(l)$ has $(r-1)$-st derivative $\asymp h\lambda$, so the
inductive bound applies to each $h$; optimising $H$ in the resulting
inequality gives the stated exponents.  (The optimisation is the
standard van der Corput $A$-process calculation.)
\end{proof}

\begin{lemma}[Type I bound]\label{lem:typeI}
Let $j\ge1$, $F>0$, and let $g(l)=F\,(L/l)^{j}$ on $l\sim L$.  For any
$r\ge2$,
\[
\Bigl|\sum_{l\sim L}e(g(l))\Bigr|
\ll_{j,r} L\bigl(F/L^{2}\bigr)^{1/(2^{r}-2)}\,
\bigl(j^{r}\bigr)^{1/(2^{r}-2)}
+L^{1-2^{2-r}}\bigl(F j^{r}/L^{r}\bigr)^{-1/(2^{r}-2)} .
\]
In particular, if $L^{\eta}\le F\le L^{r-\eta}$ for some $\eta>0$ and
$j^{r}\le L^{\eta/2}$, then
$\bigl|\sum_{l\sim L}e(g(l))\bigr|\ll L^{1-\eta'}$ for some
$\eta'=\eta'(\eta,r)>0$.
\end{lemma}

\begin{proof}
$g^{(r)}(l)\asymp_{j,r}F j^{r}/L^{r}$ on $l\sim L$, with fixed sign;
apply Lemma~\ref{lem:vdcr} with $\lambda\asymp Fj^{r}/L^{r}$ and
simplify.
\end{proof}

\section{Vaughan decomposition}\label{sec:vaughan}

\begin{lemma}[Vaughan]\label{lem:vaughan}
Let $U\ge2$ and $Q\ge U^{2}$.  For any function
$\phi:(Q,2Q]\to\mathbb C$ with $|\phi|\le1$,
\[
\Bigl|\sum_{p\sim Q}\phi(p)\Bigr|
\;\ll\;\frac{Q}{U}+
(\log Q)^{2}\max_{M\le U^{2}}\;\Bigl|\sum_{m\sim M}\;
\sum_{l:\ ml\sim Q}\phi(ml)\Bigr|
+(\log Q)^{2}\!\!\max_{\substack{U\le M\le Q/U}}\;
\bigl|\mathcal B_{M}(\phi)\bigr| ,
\]
where $\mathcal B_{M}(\phi)=\sum_{m\sim M}\sum_{l:\ ml\sim Q}
a_{m}b_{l}\,\phi(ml)$ for certain coefficients with
$|a_m|\le\tau(m)$, $|b_l|\le\log l$.
\end{lemma}

\begin{proof}
Apply the identity
$\Lambda=\mu_{\le U}*\log-\mu_{\le U}*\Lambda_{\le U}*1
+\mu_{>U}*\Lambda_{>U}*1$ to
$\sum_{q\sim Q}\Lambda(q)\phi(q)$, split all convolution ranges
dyadically, and pass from $\Lambda$-weights to primes by partial
summation (the prime powers $q=p^{i}$, $i\ge2$, contribute
$O(Q^{1/2}\log Q)$).  The first two convolutions produce the Type~I
term (the inner variable carries the smooth weight $\log$ or $1$,
removable by partial summation); the third produces the bilinear term
with the stated coefficient bounds.  Full details are the standard
Vaughan computation and are reproduced in
Appendix~A.\ \emph{[Appendix A to be written out.]}
\end{proof}

\section{The bilinear bound}\label{sec:typeII}

\begin{lemma}[Type II bound]\label{lem:typeII}
Let $j\ge1$, and let $F>0$ denote the amplitude of
$\Phi(m,l)=F\,(MQ/(ml)\cdot Q^{-1})^{j}$ (so that
$\Phi(m,l)=h\,n/(ml)^{j}$ with $F=h\,n/Q^{j}$ on $ml\sim Q$, $m\sim M$).
Suppose $Q^{1/3}\le M\le Q^{1/2}$ and, for some $\eta>0$ and $r\ge2$,
$L^{\eta}\le F\,|\Delta|/L\le L^{r-\eta}$ holds for the differenced
amplitudes ($L=Q/M$, $1\le|\Delta|\le L$).  Then for coefficients
$|a_m|\le\tau(m)$, $|b_l|\le\log l$,
\[
|\mathcal B_{M}|\;\ll\;Q\,(\log Q)^{3}
\Bigl(M^{-1/2}+L^{-\eta'/2}\Bigr)
\]
with $\eta'=\eta'(\eta,r)$ as in Lemma~\ref{lem:typeI}.
\end{lemma}

\begin{proof}
By Cauchy--Schwarz in $m$ (absorbing $\tau(m)^{2}$ by the standard mean
value $\sum_{m\sim M}\tau(m)^{2}\ll M(\log M)^{3}$),
\[
|\mathcal B_{M}|^{2}\ll M(\log Q)^{3}
\sum_{l,l'\sim L}|b_l b_{l'}|
\Bigl|\sum_{m\sim M}e\bigl(\Phi(m,l)-\Phi(m,l')\bigr)\Bigr| .
\]
The diagonal $l=l'$ contributes $\ll M L(\log Q)^{2}\cdot M
= M^{2}L(\log Q)^{2}$.  For $l\ne l'$, the differenced phase in $m$ is
$h\,n\,m^{-j}(l^{-j}-l'^{-j})$, a monomial in $m$ of amplitude
$\asymp F\,|l-l'|/L\cdot(Q/(ML))^{j}\asymp F|l-l'|/L$; apply
Lemma~\ref{lem:typeI} in the variable $m$ and sum over $l\neq l'$.
Combining and taking square roots gives the claim.
\end{proof}

\section{Scheduling: covering every regime}\label{sec:schedule}

The parameters are chosen diagonally in $n$:
\[
\varepsilon_n=(\log\log n)^{-1/3},\qquad
J(n)=\bigl\lceil(\log\log n)^{1/3}\bigr\rceil,\qquad
K(n)=(\log n)^{1-\varepsilon_n},\qquad
H(n)=(\log n)^{2} .
\]
For a band $k\le K(n)$, a dyadic block $p\sim Q$ inside it, and a level
index $i\le J(n)$, the amplitude is $F_i=h\,n/Q^{\,k-i+1}$ with
$1\le h\le H$; across the band $F_i$ ranges in
$[h,\;h\,Q^{\,i}]$, so all amplitudes obey
$F\le H\,Q^{\,J+1}$.  Taking $r=J+3$ in
Lemmas~\ref{lem:typeI} and~\ref{lem:typeII}, the saving exponent is
\[
\eta'\gg 2^{-r}=2^{-J-3}\ge\exp\bigl(-c(\log\log n)^{1/3}\bigr),
\qquad\text{so}\qquad
L^{\eta'}\ge\exp\Bigl(c'\,(\log n)^{\varepsilon_n}\,
e^{-c(\log\log n)^{1/3}}\Bigr)\;\gg\;(\log n)^{A}
\]
for every fixed $A$, uniformly over $k\le K(n)$: the block length
satisfies $\log L\gg\log Q\gg(\log n)^{\varepsilon_n}
=\exp\bigl((\log\log n)^{2/3}\bigr)$, which dominates
$2^{J}=\exp\bigl(O((\log\log n)^{1/3})\bigr)$.  \emph{The full regime
verification --- that every $(F,Q,i,k)$ triple arising in the theorem
falls into either the Kusmin--Landau range ($F\le L^{\eta}$: the phase
is nearly stationary and the interval count is controlled directly by
the band geometry of \S\ref{sec:density}) or the derivative-test range
of Lemma~\ref{lem:typeI} --- is the final audit pass of this paper and
is set out case-by-case in Appendix~B.  [Appendix B to be written
out.]}

Combining Lemmas~\ref{lem:vaughan}--\ref{lem:typeII} with the
$J$-dimensional Erd\H{o}s--Tur\'an inequality (proved with product
Fej\'er kernels in Appendix~C; the weak form with $(\log)$-power losses
suffices) yields Proposition~\ref{prop:equi} with
$o(1)\le(\log n)^{-1}$, uniformly in $n$ and $k\le K(n)$.

\section{Assembly}\label{sec:assembly}

Assuming Proposition~\ref{prop:equi}:
\[
E(n)\;\le\;
\sum_{k\le K}\Bigl(2^{-J(k)}+o(1)\Bigr)\Bigl(\log\bigl(1+\tfrac1k\bigr)+O(\cdot)\Bigr)
\;+\;\sum_{K<k}\frac{C}{k}\Big|_{k\le\log_3 n}
\]
with $K=(\log n)^{1-\varepsilon}$.  The first sum is $O(1)$ by the choice of
$J$; the second is $\le C\log\frac{\log n}{K}=C\varepsilon'\log\log n$.
Letting $\varepsilon\to0$ slowly gives $E(n)=o(\log\log n)$.
\emph{[Constants and the $O(\cdot)$ band-Mertens error to be written; the
band-Mertens statement with explicit error is already available in
machine-verified form as \texttt{primeHarmonicMass\_railBand\_le\_uniform}.]}

\section*{Provenance and verification}

The route is Tao's sketch (Erd\H{o}s problem registry no.~377, comment of
29~Aug 2025); the contribution here is its execution.  Components (K) and
(B), the conservation law
$\sum_{p}\bigl(\text{carries}\bigr)\log p=\log\binom{2n}{n}$, and the
reduction architecture are formalized in Lean~4 over Mathlib
(\texttt{RequestProject/Erdos377CarryRails.lean}, $\sim$7{,}600 lines, axiom
footprint $\{\texttt{propext},\texttt{Classical.choice},\texttt{Quot.sound}\}$,
no \texttt{sorry}).  Component (E) is classical analytic number theory and is
not formalized.

\end{document}
